\documentclass[11pt, a4paper]{article}

\usepackage[T1]{fontenc}
\usepackage[utf8]{inputenc}
\usepackage{tgpagella}
\usepackage[spanish, es-noshorthands]{babel}
\usepackage{booktabs}
\usepackage{tabularx}
\usepackage[table, xcdraw]{xcolor}
\usepackage[most]{tcolorbox}
\usepackage[margin=1.5cm]{geometry}

\definecolor{ucbblue}{RGB}{0, 51, 102}

\begin{document}

\begin{center}
    \Large\textbf{\textcolor{ucbblue}{Registro de Resultados: Recuperatorio 1 (EC1)}}[cite: 9]
\end{center}

\textbf{Curso:} IMT-221 — Circuitos Electrónicos II \hfill \textbf{Elemento de Competencia:} EC1 — Análisis Dinámico de CA[cite: 9]

\vspace{0.4cm}
\begin{tcolorbox}[colback=white, colframe=ucbblue, title=\textbf{Tabla de Seguimiento Analítico[cite: 9]}]
    \renewcommand{\arraystretch}{1.5}
    \small
    \begin{tabularx}{\textwidth}{@{} X c c c c c c c @{}}
    \toprule
    \rowcolor[HTML]{EFEFEF}
    \textbf{Estudiante} & \textbf{Nota /50} & \textbf{Modelado} & \textbf{Formulación} & \textbf{Álgebra} & \textbf{Interp.} & \textbf{Resultado EC1} & \textbf{Follow-up} \\ \midrule
    & & & & & & & \\ \hline
    & & & & & & & \\ \hline
    & & & & & & & \\ \hline
    & & & & & & & \\ \hline
    & & & & & & & \\ \hline
    & & & & & & & \\ \hline
    & & & & & & & \\ \hline
    & & & & & & & \\
    \bottomrule
    \end{tabularx}
\end{tcolorbox}

\vspace{0.4cm}
\textbf{Separación Nota-Diagnóstico:} La nota numérica sumativa y el diagnóstico por dimensiones son métricas independientes. Dos estudiantes con el mismo puntaje pueden presentar déficits operativos distintos que requieran planes de seguimiento diferentes[cite: 9].

\vspace{0.4cm}
\begin{tcolorbox}[colback=white, colframe=black, title=\textbf{Dimensiones de la Competencia (Adecuada / Parcial / Insuficiente)[cite: 9]}]
    \begin{itemize}
        \item \textbf{Modeling Competence:} ¿Transformó correctamente el circuito físico en un modelo fasorial?[cite: 9]
        \item \textbf{Equation-Formulation:} ¿Construyó relaciones coherentes desde la topología empleando KCL/KVL/divisores?[cite: 9]
        \item \textbf{Complex-Analysis:} ¿Pudo desarrollar el cálculo matemático complejo sin destruir el modelo estructural?[cite: 9]
        \item \textbf{Interpretation:} ¿Pudo convertir el resultado matemático en una afirmación eléctrica razonable y físicamente plausible?[cite: 9]
    \end{itemize}
\end{tcolorbox}

\vspace{0.4cm}
\begin{tcolorbox}[colback=gray!10, colframe=gray!50!black, title=\textbf{Etiquetas Diagnósticas (Para registrar en Follow-up)[cite: 9]}]
    \texttt{[Interpretation]} \quad \texttt{[Circuit model]} \quad \texttt{[Impedance]} \quad \texttt{[Method selection]} \\
    \texttt{[KCL/KVL]} \quad \texttt{[Sign convention]} \quad \texttt{[Complex arithmetic]} \quad \texttt{[Algebra]} \\
    \texttt{[Transfer function]} \quad \texttt{[Magnitude/phase]} \quad \texttt{[Physical interpretation]}[cite: 9]
\end{tcolorbox}

\end{document}
